\documentclass{article}

\textwidth=6.5in
\textheight=8.5in
\voffset=-54pt
\oddsidemargin=0in
\evensidemargin=0in
\setlength{\parskip}{8pt}
\setlength{\parindent}{0pt}

\usepackage{amsmath, amssymb}
\usepackage{ifthen}

\usepackage[inline]{enumitem}
\setlist{topsep=0pt, parsep=8pt}

\usepackage[rgb,table]{xcolor}\convertcolorsUfalse
\usepackage{graphicx}

% change hyperref options
\PassOptionsToPackage{hyphens}{url}
\usepackage[bookmarks,colorlinks,linkcolor=black,citecolor=blue,pdfstartview=FitH,urlcolor=blue]{hyperref}
\hypersetup{pdfpagemode=UseNone}
\usepackage{bookmark}

\usepackage{graphicx}
\usepackage{multicol}
\usepackage{framed}
\usepackage{array}

\usepackage{icomma}

\usepackage{marginnote}
\usepackage[colorinlistoftodos]{todonotes}
\renewcommand{\marginpar}{\marginnote}

\usepackage{soul}
\usepackage[normalem]{ulem}

\usepackage{pgfplots}
\pgfplotsset{compat=1.18}  
\pgfplotscreateplotcyclelist{mycolor}{black, blue, red, green!70!black, magenta, purple, cyan, orange }  
\pgfplotsset{
  disabledatascaling,
  cycle list=tikzcolor, 
  axis lines=middle,
  every axis plot/.style={no marks, thick},
  every axis/.style={
    enlarge x limits=false,
    enlarge y limits=auto,
    xlabel style={at={(current axis.right of origin)},anchor=west},
    ylabel style={at={(current axis.above origin)},anchor=south},
    % extra x and y ticks for asymptotes
    extra x tick style={grid=major, grid style={dashed,black}},
    extra y tick style={grid=major, grid style={dashed,black}},
    /pgf/number format/1000 sep={},
  },    
  tick label style = {font=\footnotesize, fill=\currentbgcolor, inner sep=0pt},    
  xticklabel shift=1pt,    
  scaled ticks=false,
  scaled y ticks=false,
  unbounded coords=jump,
  samples=101,
  minor tick length=0,
  scatter/classes={
    open={mark=*,draw=black,fill=white, mark size=2, thin},
    closed={mark=*,draw=black,fill=black, mark size=2, thin}
  },
  width=3in,
}
\pgfplotsset{open/.style={only marks, mark=*, fill=white, thin, mark size=2}}
\pgfplotsset{closed/.style={only marks, mark=*, draw=black, fill=black,
    thin, mark size=2}}
\pgfplotsset{labeled/.style={nodes near coords, only marks, mark=*, draw=black, fill=black, thin, mark size=2, point meta=explicit symbolic}}

\usepackage{tikz}
\usetikzlibrary{
  matrix,%
  % enables the snake paths
  decorations.pathmorphing,%
  % enables bigger arrows
  decorations.markings,%
  decorations.pathreplacing,%
  decorations.text,%
  calc,%
  patterns,%
  shapes.geometric,%
  arrows,
  decorations.shapes,
  positioning,
  plotmarks,
  intersections
}
\tikzset{>=latex} % sets default arrow type in tikz
\tikzset{font=\normalsize}
\tikzset{mark size=1.5pt, mark options=thin}
\tikzset{pin distance=4pt,
  every pin edge/.style={<-, thin, shorten <= -2pt}}

\interfootnotelinepenalty=10000
\renewcommand{\thefootnote}{(\arabic{footnote})}

\usepackage{multirow, bigdelim}

% change to \usepackage[sols]{handout} to see solutions
\usepackage[sols, notes, grading]{handout}

\newcommand\iscurrentenumi[1]{%
  \ifthenelse{\equal{#1}{\theenumi}}{}{#1}}
\setenumerate[1]{ref=\#\arabic*}
\setenumerate[2]{ref=\protect\iscurrentenumi{\theenumi}(\alph*)}
\newcommand\enumiiprefix[2]{%
  \ifthenelse{{\equal{#1}{\theenumi}}\AND{\equal{#2}{(\alph{enumii})}}}{}{}%
  \ifthenelse{{\equal{#1}{\theenumi}}\AND{\NOT{\equal{#2}{(\alph{enumii})}}}}{#2}{}%
  \ifthenelse{\equal{#1}{\theenumi}}{}{#1#2}%
}
\setenumerate[3]{ref=\protect\enumiiprefix{\theenumi}{(\alph{enumii})}\roman*}

\makeatletter

% underline links               
\let\@href=\href
\renewcommand{\href}[2]{\@href{#1}{\ul{#2}}}

\newcommand\calculator{
  \begin{tikzpicture}[x=0.15ex, y=0.15ex, baseline=1.5pt]
    \begin{scope}[thick]
      \draw (0, 0) -- (10, 0) -- (10, 14) -- (0, 14) -- cycle;
      \foreach \x in {10/3, 20/3}
      {
        \draw (\x, 0) -- (\x, 10);
      }
      \foreach \y in {10/3, 20/3, 10}
      {
        \draw (0, \y) -- (10, \y);
      }
    \end{scope}
  \end{tikzpicture}
}
\newcommand{\calcitem}{\refstepcounter{enum\romannumeral\@enumdepth}%
\item[\calculator \csname labelenum\romannumeral\@enumdepth\endcsname]}

\makeatother

\definecolor{purple}{rgb}{0.5,0,1.0}

\newcommand{\doedfinity}[2][\newpage]{
  \begin{problem}[#1]
    Do the Edfinity assignment ``Problem Set #2''.

    We encourage you to write down any scratch work here, as well as
    things you want your future self to remember. Nothing you write
    here will be graded; it's just for you!
  \end{problem}
}

\newcommand{\psrnewpage}{\newpage}
\newcommand{\psreflection}[2][]{

  \ifthenelse{\not\equal{#1}{nocite}}{
    \begin{enumerate}[resume]
    \item
      \begin{problem}[1in]
        Please cite any help you got on this problem set:
        \begin{itemize}[nosep]
        \item If you worked with other students, please list their names here.
        \item If you used resources other than those on our Canvas site, please list them here.
        \end{itemize}
      \end{problem}

    \end{enumerate}
  }

  \probonly{%
    #2

    \psrnewpage

    \emph{Reflection space.} It's a good habit to take a few minutes at the end of each pset to reflect on what you learned and jot down notes to your future self. Here are a few reflection questions to get you started. Nothing you write here will be graded; it's just for you!
    \begin{itemize}
    \item Take a few minutes to look back at the problems. How could you explain to someone else how to do each problem? What was your strategy, and how did you know to use that strategy? If there are problems where you don't feel able to answer this, write down which ones they are. Come back to them in a few days when we post the homework solutions!

      \vspace{1.5in}

    \item Take a look back at the learning objectives on today's worksheet; how are you feeling about each one? If there are ones you know you need more practice with, write those down here.

      \vspace{1.5in}

    \item What did you find most challenging about this material? Do you have any lingering questions about it? If so, write them down here so you don't forget them. At some point, stop by office hours or MQC to ask!

      \vfill
    \end{itemize}      
  }
}

\usepackage{fancyhdr}
\lhead{\textcolor{white}{This content is protected and may not be shared, uploaded, or distributed.}}
\rhead{}
\rfoot{\vspace*{\fill}\footnotesize\textcopyright{} \emph{Harvard Math Ma}}
\renewcommand{\headrulewidth}{0pt}
\pagestyle{fancy}
\begin{document}

\begin{center}
  \large\textsc{Problem Set 26 - Modeling with Exponential Functions}\vspace{.5in}\\
  \begin{problemonly}
   \large Name: \rule{5cm}{0.15mm}\\
   \end{problemonly}
\end{center}

\begin{enumerate}
\item  \note{
    \begin{itemize}[nosep]
    \item Some limits involving exponentials
    \item Derivative of something like $\sqrt[3]{\pi e^x}$
    \end{itemize}
  }

  \doedfinity{26}

\item 

  In June 2023, former basketball star Michael Jordan decided to sell his stake in the Charlotte Hornets, a basketball team. \href{https://www.nytimes.com/2023/06/16/sports/basketball/michael-jordan-charlotte-hornets-sale.html?unlocked_article_code=1.WE4.efbp.a6g0fLAio0T-&smid=url-share}{According to the New York Times}, the team was valued at \$3 billion in the deal. Jordan bought his stake in 2010, when the team was valued at \$275 million.

  Let's find the average annual percent change in the team's valuation from 2010 to 2023. Here's how we do that:
  \begin{enumerate}
  \item

    \begin{grading}
      2 points: 0.5 for having a formula with $v(0) = 275$, 1.5 for the rest of the formula
    \end{grading}

    \begin{problem}[\vfill]
      Assuming the valuation is an exponential function of time, write a formula for $v(t)$, the team's valuation in \textbf{millions} of dollars $t$ years after 2010. 
    \end{problem}

    \begin{solution}
      We know that $v(0)=275$ and $v(13)=3000$. To get from $275$ to $3000$, we need
      to multiply by $\frac{3000}{275}=\frac{120}{11}$. Since this happens over a period
      of 13 years, $v(t) = 275\left(\frac{120}{11}\right)^{t/13}$.
    \end{solution}

  \calcitem

    \begin{grading}
      1 point
    \end{grading}

    \begin{problem}[\vfill]
      By what percent does $v(t)$ increase each year? This is called the \emph{average annual percent change}. 
    \end{problem}

    \begin{solution}
      Each year, the amount is multiplied by $\left(\frac{120}{11}\right)^{1/13}$, so
      the average annual percent change is 
      $\left(\frac{120}{11}\right)^{1/13}-1\approx 0.202$, or about
      $20.2\%$.
    \end{solution}

    \probonly{For comparison, a person who invests in regular stocks for many years can expect the investment to earn an average of about 10\% per year.} 
  \end{enumerate}

  \pspace{\newpage}

\item 

  Throughout the COVID pandemic, epidemiologists around the world tried to model the spread of the disease. Here is \href{https://www.desmos.com/calculator/x400n2fvrh}{data showing the number of daily new confirmed COVID cases between March 1, 2020 and March 25, 2020.} The independent variable is time in days, with $t = 0$ representing March 1, 2020; the dependent variable is the number of new confirmed COVID cases that day.

  \begin{enumerate}
  \item

    \begin{grading}
      2 points: 0.5 for having a formula that goes through $(0, 1793)$, 1.5 for having the rest of the formula
    \end{grading}

    \begin{problem}[\vfill]
      Looking at the graph, we can see right away that it doesn't look linear. Let's try modeling it with an exponential function. Find the exponential function that goes through the first and last data points, $(0, 1793)$ and $(24, 47226)$.
    \end{problem}

    \begin{solution}
      $f(t) = 1793\left(\frac{47226}{1793}\right)^{t/24}$.
    \end{solution}

  \item

    \begin{grading}
      0.5 points. (A yes/no answer is fine.) 
    \end{grading}

    \begin{problem}[\vfill\newpage]
      Plot your exponential function together with the data. (Scroll down past the table, and you'll see a line where you can enter your function.) Does your exponential function seem like a good match for the data?
    \end{problem}

    \begin{solution}
      Yes, the exponential function fits the data well.
    \end{solution}

  \item

    \begin{grading}
      1.5 points: the key is that we're assuming that, on average, each person infects the same number of other people 
    \end{grading}

    \begin{problem}[\vfill]
      When we model data using a particular type of function (like a linear function or an exponential function), we would really like to have some reason (beyond just the data) to think that we've picked an appropriate type of function.

      Read the first paragraph of \href{https://canvas.harvard.edu/files/21213978/download?download_frd=1}{this article}, and explain why this suggests that an exponential model makes sense for the spread of COVID in this time period. 
    \end{problem}

    \begin{solution}
      If we assume that each infected person infects the same number of additional 
      people,
      then this results in a constant percent change, which is best modeled using
      an exponential model.
    \end{solution}

  \item

    \begin{grading}
      2 points: 0.5 for saying it isn't reasonable to model it using an exponential function, 1.5 for any thoughtful discussion of real world factors
    \end{grading}

    \begin{problem}[\vfill\newpage]
      \href{https://ourworldindata.org/explorers/coronavirus-data-explorer?zoomToSelection=true&time=2020-03-01..2020-05-01&facet=none&country=~OWID_WRL&pickerSort=asc&pickerMetric=location&Interval=7-day+rolling+average&Relative+to+Population=false&Color+by+test+positivity=false&Metric=Confirmed+cases}{This graph} shows the data for COVID cases between March 1, 2020 to May 1, 2020.

      Does it seem reasonable to model this using an exponential function? What real world factors make this situation different from the situation between March 1, 2020 and March 25, 2020? 
    \end{problem}

    \begin{solution}
      It does not seem reasonable to model the COVID cases in this period using an exponential
      function, because the rate of change of the number of new cases decreases drastically
      in April. This may have been due to social distancing and lockdown policies.
    \end{solution}

  \end{enumerate}

\item %({\it 9.2.26-29})

  \begin{enumerate}
  \item In each part, you're given a graph. Find a function of the form $f(x) = C a^x + D$ that fits the graph. Show your work. 

    {\it Hint:} You should be able to determine $D$ just by looking at the graph; why? Once you've found $D$, you can shift the graph to see what $y = Ca^x$ looks like.

    \begin{grading} 
      3.5 points per part:
      \begin{itemize}[nosep]
      \item 0.5 for correctly identifying $D$
      \item 1 for correctly setting up a system of equations about $a$ and $C$
      \item 1 for correctly solving the system for $a$
      \item 1 for correctly using that to find $C$
      \end{itemize}
    \end{grading}

    \probonly{\begin{multicols}{2}}
      \begin{enumerate}

      \item
        \begin{problem}
          \raisebox{2ex-\height}{
            \begin{tikzpicture}
              \begin{axis}[xtick=\empty, ytick=\empty, clip=false, width=2.25in]
                \addplot+[domain=-4:2.5] {8*(2)^(x)-32}; %
                \addplot+[domain=-4:2.5, dashed] {-32} node[right]{$y = -32$}; %
                \addplot+[closed] coordinates {(-3, -31)} node[anchor=south east, inner sep=1pt]{$(-3, -31)$}; %
                \addplot+[closed] coordinates {(1, -16)} node[anchor=north west, inner sep=1pt]{$(1, -16)$};%
              \end{axis}
            \end{tikzpicture}
          }
        \end{problem}

        \begin{solution}
          Since the horizontal asymptote is $y=-32$, $D=-32$. Using the points
          $(-3, 31)$ and $(1, -16)$, we can write down a system of equations
          to solve for $C$ and $a$.
          \begin{align*}
            Ca^{-3} - 32 &= -31 \\
            Ca^1 - 32 &= -16
            \intertext{Adding 32 to both sides of both equations,}
            Ca^{-3} &= 1 \\
            Ca &= 16
          \end{align*}
          Now we can divide the second equation by the first equation to get rid of $C$.
          \begin{align*}
            \frac{Ca}{Ca^{-3}} &= \frac{16}{1} \\
            \frac{a}{a^{-3}} &= 16 \\
            a^4 &= 16\\
            a &= 2
          \end{align*}
          We can then plug $a=2$ into the equation $Ca=16$ to find that $C=8$.
          
          Thus, $\boxed{f(x)=8(2)^x-32}$.
        \end{solution}

      \item

        \begin{problem}
          \raisebox{2ex-\height}{
            \begin{tikzpicture}
              \begin{axis}[xtick=\empty, ytick=\empty, clip=false, width=2.25in]
                \addplot+[domain=-5:3] {6-(1/2)^((x-1)/2)}; %
                \addplot+[domain=-5:3, dashed] {6} node[right]{$y = 6$}; %
                \addplot+[closed] coordinates {(1, 5)} node[anchor=north west, inner sep=1pt]{$(1, 5)$};%
                \addplot+[closed] coordinates {(-1, 4)} node[anchor=south east, inner sep=1pt]{$(-1, 4)$};
              \end{axis}
            \end{tikzpicture}
          }
        \end{problem}

        \begin{solution}
          Since the horizontal asymptote is $y=6$, $D=6$. Using the points
          $(-1, 4)$ and $(1, 5)$, we can write down a system of equations
          to solve for $C$ and $a$.
          \begin{align*}
            Ca^{-1} + 6 &= 4 \\
            Ca^1 + 6 &= 5
            \intertext{Subtracting 6 from both sides of both equations,}
            Ca^{-1} &= -2 \\
            Ca &= -1
          \end{align*}
          Now we can divide the second equation by the first equation to get rid of $C$.
          \begin{align*}
            \frac{Ca}{Ca^{-1}} &= \frac{-1}{-2} \\
            \frac{a}{a^{-1}} &= \tfrac{1}{2} \\
            a^2 &= \tfrac{1}{2}\\
            a &= \left(\tfrac{1}{2}\right)^{1/2}
          \end{align*}
          We can then plug $a=\left(\tfrac{1}{2}\right)^{1/2} $ into the equation 
          $Ca=-1$ to find that $C=-2^{1/2}$.
          
          Thus, $\boxed{f(x)=-2^{1/2}\left(\tfrac{1}{2}\right)^{x/2}+6}$.
        \end{solution}

      \end{enumerate}
    \probonly{\end{multicols}}

  \pspace{\vfill}

  \item

    \begin{grading}
      2 points: 1 for finding the $x$-intercept of (i) correctly, and 1 for a reasonable explanation of why we can't find the other by hand. 
    \end{grading}

    \begin{problem}[\nstretch{0.5}]
      You can find the exact decimal value (by hand) of the $x$-intercept of \uline{one} of the graphs in (a). Which one? Find the $x$-intercept for that graph, and explain why you can't find the other by hand. 
    \end{problem}

    \begin{solution}
      To find the $x$-intercept of either graph, we need to solve the equation $f(x)=0$.
      For the first graph,
      \begin{align*}
        8(2)^x-32 &= 0 \\
        8(2)^x &= 32 \\
        2^x &= 4 \\
        x &= 2
      \end{align*}
      For the second graph,
      \begin{align*}
          -2^{1/2}\left(\tfrac{1}{2}\right)^{x/2} + 6 &= 0 \\
          2^{1/2}\left(\tfrac{1}{2}\right)^{x/2} &= 6 \\
          \left(\tfrac{1}{2}\right)^{x/2} &= \frac{6}{2^{1/2}} \\
          2^{-x/2} &= \frac{6}{2^{1/2}}
      \end{align*}
      We cannot solve this equation by hand because $6$ cannot be written as
      a simple power of $2$.
    \end{solution}
  \end{enumerate}

\item %(9.4.14)

  Money in a bank account is growing exponentially.  When there's \$4000 in the account, the account is growing at an instantaneous rate of \$100 per year.  

  \begin{enumerate}
  \item

    \begin{grading}
      1 point
    \end{grading}

    \begin{problem}[\vfill]
      How quickly is the money growing when there's \$5000 in the account? Explain how you know. 
    \end{problem}

    \begin{solution}
      Since $M(t)$ is an exponential fucntion, we know that 
      $M'(t) = kM(t)$, where $k$ is a constant of proportionality.
      When $M(t)=4000$, $M'(t)=100$, which means $k=\frac{100}{4000}=\frac{1}{40}$.

      So when $M(t)=5000$, $P'(t)=\frac{1}{40}(5000)=125$. Therefore
      the money is growing at \fbox{\$125/year} when there is \$5,000 in the account.
    \end{solution}

  \item

    \begin{grading}
      1 point
    \end{grading}

    \begin{problem}[\vfill]
      We know that exponential functions change at a rate proportional to their size. What's the constant of proportionality in this case? In other words, fill in the blank: $M'(t) = \probonly{\vphantom{\dfrac{1}{1}}}\cfitb{0.5in}{\frac{1}{40}} M(t)$.
    \end{problem}

    \begin{solution}
      We found in part (a) that the constant of proportionality is \fbox{$\frac{1}{40}$}.
    \end{solution}

  \item

    \begin{grading}
      2 points
    \end{grading}

    \begin{problem}[\stretch{1.25}]
      Let $t$ be measured in years, and suppose that $M(0) = 3000$. One of the following is a formula for $M(t)$; which one? Explain how you know. 
      \begin{multicols}{4}
        \begin{enumerate}
        \item $3000 e^{0.025 t}$
        \item $3000 e^{40 t}$
        \item $4000 e^{0.025 t}$
        \item $4000 e^{40 t}$
        \end{enumerate}
      \end{multicols}
    \end{problem}

    \begin{solution}
      We know that $M'(t)=\frac{1}{40}M(t)=0.025M(t)$, so that rules out (ii) and (iv).
      Since we know that $M(0)=3000$, that rules out (iii), leaving us with \fbox{(i)}. 
    \end{solution}

  \item

    \begin{grading}
      1 point
    \end{grading}

    \begin{problem}[\nstretch{0.5}]
      Since $M(t)$ is an exponential function, it can be written as $C \cdot b^t$ for some constants $C$ and $b$. What are the values of $C$ and $b$? 
    \end{problem}

    \begin{solution}
      $M(t)=3000e^{0.025t}$. Based on this, \fbox{$C=3000$ and $b=e^{0.025}$}.
    \end{solution}

  \end{enumerate}
\item  Exam 2 is on Thursday 11/14. Please read \href{https://canvas.harvard.edu/courses/138327/pages/exams}{the information on Canvas about the exam}.

  \begin{enumerate}
  \item

    \begin{grading}
      1.5 points, full credit for any thoughtful response
    \end{grading}

    \begin{problem}[\stretch{0.25}]
      After Exam 1, you filled out a learning reflection survey; \href{https://forms.gle/gZoESAGCsZBYJ3VC9}{reread your response now}. You likely identified opportunities to study more effectively. How will you incorporate those into your study plan for Exam 2? What's one thing you plan to do as a result of your experience taking Exam 1?
    \end{problem}

  \item

    \begin{grading}
      1.5 points, full credit for any thoughtful response
    \end{grading}

    \begin{problem}[\vfill\newpage]
      As you did in \href{https://canvas.harvard.edu/courses/138327/files/20808062?wrap=1}{Problem Set 11}, {{\#6}}, work backwards from our exam date (Thursday 11/14) to make a schedule for your Exam 2 study activities. As you do this, keep in mind what you've learned in workshop about testing yourself and spacing out your study sessions!

      We've posted 3 practice exams, and we strongly encourage you to do all three all the way through (and we don't mean just reading the solutions!). 
    \end{problem}

  \end{enumerate}

\end{enumerate}
\psreflection{}
\end{document}
